\documentclass[10pt,a4paper]{amsart}
\usepackage[margin=19mm]{geometry}
\usepackage{amsmath,amssymb,mathtools}
\usepackage[scr=rsfs,cal=euler]{mathalfa}
\usepackage{xfrac}
\usepackage[unicode,hidelinks,bookmarks=false]{hyperref}
\newcommand{\ie}{i.e.\ }
\newcommand{\cf}{cf.\ }
\newcommand{\resp}{respectively\ }
\newcommand{\Subsection}{\subsection}
\providecommand{\nobreakdash}{\nobreak\hbox{-}}
\setlength{\emergencystretch}{2em}
\multlinegap0pt
\usepackage{amssymb}
\usepackage{mathtools}
\usepackage[shortlabels]{enumitem}
\newtheorem{corollary}{Corollary}
\newtheorem{lemma}{Lemma}
\newcommand{\R}{\mathbb R}
\DeclareMathOperator{\pos}{pos}
\title{The Discrete $L_p$ Minkowski Problem for Negative $p$}
\author{Junjie Shan}
\date{Source: arXiv:2609.01429v1 (CC BY 4.0)}
\begin{document}
\maketitle

\noindent\emph{Source and licence.} The Discrete $L_p$ Minkowski Problem for Negative $p$. Version arXiv:2609.01429v1 (CC BY 4.0), distributed by arXiv under the Creative Commons Attribution 4.0 International licence, http://creativecommons.org/licenses/by/4.0/, as recorded in the arXiv licence metadata for this exact version and shown on the abstract page https://arxiv.org/abs/2609.01429v1. Full licence notice: \texttt{licenses/arxiv-2609.01429-CC-BY-4.0.txt}.\par\smallskip

\noindent\textit{Editorial scope.} Counted unit: the article's lemma lem:W-regularity (source0.tex lines 769--779). The article's complete original proof of it is delivered with it (source0.tex lines 781--909, one \texttt{proof} environment of 129 lines). No mathematics is written, repaired or re-derived: the statement and the proof are the article's own lines, in the article's own notation, and no retained line is substituted or re-flowed.  Retained uncounted as the source-local context the proof needs: lines 301--308; lines 652--662.  Nothing was omitted from any retained span, so the package carries no omission record.  No retained line contains a citation, so the wrapper carries no bibliography and no bibliography entry.  No retained line contains a figure environment, an image-inclusion command or drawing code, so no image is delivered and the package declares no asset dependency.  Editorial formatting only, in addition: an AMS \texttt{amsart} preamble that restates the article's own declarations and macros used by the retained lines, namely the environments \texttt{corollary}, \texttt{lemma} on separate counters, together with the standard packages those lines need.  The retained environments print in this document as Corollary 1, Lemma 1 in order of appearance, whereas the article prints the same items with the numbering of its own full document.  The complete native source of the article as distributed in the arXiv e-print for this exact version is bundled unchanged as \texttt{sources/209109/source0.tex}.
\begin{equation}\label{eq:positive-span-separation}
 \begin{aligned}
  \pos U=\R^n
  &\quad\Longleftrightarrow\quad (\pos U)^\circ=\{0\}\\
  &\quad\Longleftrightarrow\quad
  U\text{ is not contained in a closed hemisphere}.
 \end{aligned}
\end{equation}

\begin{corollary}\label{cor:inverse-volume-derivative}
The function $W$ belongs to $C^1((0,\infty)^N)$ and satisfies
\begin{equation}\label{eq:inverse-volume-derivative}
 \frac{\partial W}{\partial t_i}(t)=\frac1p m_i(t).
\end{equation}
Moreover,
\begin{equation}\label{eq:inverse-volume-homogeneity}
 W(ct)=c^{n/p}W(t),
 \qquad c>0.
\end{equation}
\end{corollary}

\begin{lemma}\label{lem:W-regularity}
	The function \(W:X\to(0,\infty)\) is continuous and is
	\(C^1\) on \(\operatorname{int}\Delta_\alpha\). Moreover, for every
	\(t^0\in S\), there exist a neighborhood \(\mathcal U\) of \(t^0\)
	in \(\operatorname{aff}\Delta_\alpha\) and a function
	\(\widetilde W\in C^1(\mathcal U)\) such that
	\(
	\widetilde W=W\)
	on \(\mathcal U\cap X
	\).
\end{lemma}

\begin{proof}
	Since
	\(\operatorname{int}\Delta_\alpha\subset(0,\infty)^N\), the
	\(C^1\) assertion in the interior follows from
	Corollary~\ref{cor:inverse-volume-derivative}.
	
	Fix \(t^0\in S\), and set \(J=J(t^0)\). Since \(t^0\in X\), the
	directions \(u_j\), \(j\in J\), positively span \(\R^n\). 
	Since \(t_j^0>0\) for every \(j\in J\), set
	\(
	\eta=\frac12\min_{j\in J}t_j^0>0.
	\)
	Choose a sufficiently small neighborhood \(\mathcal U\) of \(t^0\)
	in \(\operatorname{aff}\Delta_\alpha\) such that
	\begin{equation}\label{low bound}
	t_j\geq\eta
	\qquad
	\text{for every }j\in J\text{ and }t\in\mathcal U.
	\end{equation}
	For \(t\in\mathcal U\), define
	\[
	P_J(t)
	=
	\bigcap_{j\in J}
	\{x\in\R^n:\langle x,u_j\rangle\leq t_j^{1/p}\},
	\qquad
	\widetilde W(t)=V(P_J(t)).
	\]
	Since
	\(
	\pos\{u_j:j\in J\}=\R^n,
	\)
	Corollary~\ref{cor:inverse-volume-derivative}, with the directions
	\(u_j\), \(j\in J\), in place of \(u_1,\ldots,u_N\), shows that
	\[
	(s_j)_{j\in J}
	\longmapsto
	V\left(
	\bigcap_{j\in J}
	\{x:\langle x,u_j\rangle\leq s_j^{1/p}\}
	\right)
	\]
	is \(C^1\) whenever \(s_j>0\) for every \(j\in J\).
	Since \(t_j\geq\eta>0\) on \(\mathcal U\) for \(j\in J\), and
	\(\widetilde W\) depends only on these coordinates, it follows that
	\(\widetilde W\in C^1(\mathcal U)\).
	
	
	
We next show that the polytopes \(P_J(t)\), \(t\in\mathcal U\),
are uniformly bounded. Set
	\[
	c_J
	=
	\min_{v\in S^{n-1}}
	\max_{j\in J}\langle v,u_j\rangle.
	\]
	Since
	\(\pos\{u_j:j\in J\}=\R^n\), 
	\eqref{eq:positive-span-separation} implies that, for every
	\(v\in S^{n-1}\), there exists \(j\in J\) such that
	\(\langle v,u_j\rangle>0\). Hence the continuous function
	\(
	v\longmapsto \max_{j\in J}\langle v,u_j\rangle
	\)
	is strictly positive on  \(S^{n-1}\). Therefore
	\(
	c_J
	>0
	\).
	 If
	\(x\in P_J(t)\) and \(x\neq0\), then
	\[
	c_J|x|
	\leq
	\max_{j\in J}\langle x,u_j\rangle
	\leq
	\max_{j\in J}t_j^{1/p}
	\leq
	\eta^{1/p},
	\]
	where the last inequality uses \(p<0\) and \eqref{low bound}. Hence
	\[
	P_J(t)\subset RB,
	\qquad
	R=\frac{\eta^{1/p}}{c_J},
	\qquad
	t\in\mathcal U.
	\]
	
	For every \(i\notin J\), we have \(t_i^0=0\). Since
	\(t_i^{1/p}\to+\infty\) as \(t_i\to0+\), we may shrink
	\(\mathcal U\) once more so that, whenever
	\(t\in\mathcal U\cap X\) and \(t_i>0\),
	\[
	t_i^{1/p}>R.
	\]
	Hence \(RB\subset\{x:\langle x,u_i\rangle\leq t_i^{1/p}\}\).
	Since \(P_J(t)\subset RB\), every point of \(P_J(t)\) already satisfies
	the constraint \(\langle x,u_i\rangle\leq t_i^{1/p}\) whenever
	\(i\notin J\) and \(t_i>0\). If \(t_i=0\), this constraint is absent by
	the definition of \(P(t)\). For \(t\in\mathcal U\cap X\), by \eqref{low bound} we have \(J\subset J(t)\). Hence, by the definition of
	\(P(t)\),
	\[
	P(t)
	=
	P_J(t)\cap
	\bigcap_{i\in J(t)\setminus J}
	\{x:\langle x,u_i\rangle\leq t_i^{1/p}\}.
	\]
	For every \(i\in J(t)\setminus J\), the preceding estimate shows that
	each \(x\in P_J(t)\) satisfies
	\(\langle x,u_i\rangle\leq t_i^{1/p}\). Therefore
	\[
	P(t)=P_J(t)
	\qquad
	\text{for every }t\in\mathcal U\cap X.
	\]
	
Thus, \(\widetilde W=W\) on \(\mathcal U\cap X\). Hence
\(\widetilde W\) is the required local \(C^1\) extension, and \(W\)
is continuous at \(t^0\) relative to \(X\). Since \(t^0\in S\)
was arbitrary and \(W\) is \(C^1\) on
\(\operatorname{int}\Delta_\alpha\), the equality
\(X=\operatorname{int}\Delta_\alpha\cup S\) shows that \(W\) is
continuous on \(X\).   
	

\end{proof}

\end{document}
